\chapter{A natural weak-event entry into the complete uniform C1 theorem}
\label{component:pa133fae35e52}
\noindent\textbf{Component }\path{pa133fae35e52}\par
{\small\noindent Source: \path{CORE/prior/c1/entry_and_remainder.tex}\par}
{\small\noindent SHA256: \path{a133fae35e5258c0730ee18f58ffefe20e959ddba3c6fef3905847d291e63593}\par}
{\small\noindent Inventory use: CE1--CE7 only | C1\_\allowbreak{}BARE\_\allowbreak{}OTHER\_\allowbreak{}TERMINALS; precise subordinate uses retained below\par}
\medskip
\section{A natural weak-event entry into the complete uniform C1 theorem}

Throughout, $d$ is the original positive signed coefficient. The event
coefficient at the fourth cycle position is denoted $A_{\rm ev}$, and
$D_{\rm ev}=e_5$. These must not be confused with actual state coordinates.
The source has the canonical omissions $\lambda=n=0$ and is normalized
by $c=e_3=2$. We use the canonical actual-contact sign entry, including
its admitted single-zero $D_{\rm ev}$ or $\beta$ boundaries. Explicitly,
the retained $m,c,A_{\rm ev},\alpha$ are positive, and $D_{\rm ev},\beta$
are positive or exactly one of them is zero. The noncoercive further-zero
faces are already separate exits. The inherited sign and nonsingularity
input is \path{C1_GLOBAL_SIGN_ENTRY.md}, Sections 6--7, equations
(9)--(10), on every finite noncommuting physical contact in precisely this
coefficient domain; that input is not reproved here.

\paragraph{Exact transport from the actual contact.}
Write its positive-angle contact as
\[
 u_c=k,\quad v_c=\ell,\qquad
 \psi_c=(S,k(\theta-1),kSr,1),\quad
 S>0,\quad 0<r,\theta<1.
\]
The exact C1 reflection lets us orient $0<k\le1$. Define source-only
parameters
\[
 \eta=\frac{\ell}{kS},\quad \zeta=\frac{S\ell}{k},\quad
 f=\sqrt\eta,\quad g=\sqrt\zeta,\quad
 e=\frac{k}{S}=\frac{kf}{g},\quad
 a=\frac{1-r}{\eta},\quad t=\frac\theta\zeta.
                                                               \tag{CE1}
\]
Thus the classifier $\mathfrak T=k/(S\ell)$ is exactly $\zeta^{-1}$,
and its old state-ratio variable $B$ is exactly $\eta^{-1/2}$.
Put $\kappa=k^2$, $K=1+\kappa$. The original amplitude signs imply
\[
 R=1+\kappa r,\quad T=1+\kappa(1-\theta),\quad
 M=R-t>0,\quad J=T-a>0,
\]
\[
 E=1-\kappa\eta(1-\theta)>0,\qquad
 Q_c=1-\kappa\zeta r>0.                              \tag{CE2}
\]
Indeed the positive second amplitude is $k\zeta M$, the negative
fifth amplitude is $-\ell J$, the state gap is $SE$, and the remaining
state gap is $Q_c$. In particular $0<a,t<2$, without any shape-unit
hypothesis.

The seven native coquotients used by the uniform theorem are valid
under (CE1), not just asymptotically. Their constructed contact is
\[
 (W,X,Y,Z)=(\Delta k,-k(1-t\zeta),\Delta k^2(1-a\eta),1),
 \qquad \Delta=e^{-1}=S/k,
\]
which is exactly the original contact. The native contact identities
prove all six kernel and angular equations. To make the nonsingularity
premise explicit, put $Y_0=Sr$, $h_0=S+k\ell(\theta-1)$,
$R_2=\ell(S+k^2Y_0)-k\theta$, and
$R_5=k(S-Y_0)+k^2\ell(\theta-1)-\ell$. The inherited selected
six-column contact system, on $(d,m,A_{\rm ev},D_{\rm ev},\alpha,\beta)$
with $c=2$, has determinant
\[
 -8192\theta k^2\ell(1+k^2)^2(1+\ell^2)^2\mathcal D,
\]
\[
 \mathcal D=Y_0(1+k^2)(kY_0+\ell)(1-k\ell Y_0)h_0R_2(-R_5)>0.
\]
Each factor is strictly positive by the stated source sign theorem.
Thus the normalized coefficients equal the unique original source
coefficients in that inherited system. This
establishes full original-coefficient transport, including the exact
zero openings; it does not use an archived completion label.

\paragraph{A bounded polynomial consequence of the original fifth event.}
Assume now
\[
                 k\text{ infinitesimal},\qquad
                 \zeta=\mathfrak T^{-1}\text{ infinitesimal}.     \tag{CE3}
\]
Set $q=\kappa\eta$. Since $t<2$ and $E=1-q(1-t\zeta)>0$, $q$ is
bounded and $\bar q\le1$. The original fifth-event formula is
\[
 D_{\rm ev}=\frac{a}{Jc_3}(2c_5-J\eta c_d).
\]
Using $L_c=r$, $H_c=r+\eta$ and $V_c=q\zeta$, its sign is precisely
\[
 \mathcal B_5:=K(1+q\zeta)(r+\eta)-J\eta P_d\ge0.       \tag{CE4}
\]
Although $\eta$ need not be bounded, this whole expression has the
following bounded polynomial form. Define
\begin{align*}
 H={}&2+\kappa r+\zeta[K-(\kappa+2)t]
     +q\zeta(-\zeta t^2+t+1)
     +\kappa\zeta(1-r)(t-1),\\
 P_d={}&1+q\zeta+\kappa H.
\end{align*}
Then exact cancellation using $a\eta=1-r$ and $q=\kappa\eta$ gives
\[
 \mathcal B_5=Kr(1+q\zeta)+(1-r)P_d
       +q\{\zeta t(1+q\zeta)-(K-\kappa\zeta t)H\}.     \tag{CE5}
\]
Every variable on the right is bounded. Hence (CE4) gives
\[
       1-2\bar q\ge0,\qquad 0\le\bar q\le\tfrac12.     \tag{CE6}
\]
This is obtained before dividing by $1-q$ or by a coefficient opening.
In particular $E$ is a positive unit with residue at least $1/2$.

The fourth coquotient polynomial likewise has the bounded expression
\[
\begin{split}
 P_4={}&K(1-q)+\kappa(1-r)-\kappa\zeta t\\
 &+q\zeta(-\kappa\zeta t^2+2\kappa t+t)
   +\kappa\zeta(1-r)(1-t).
\end{split}                                                       \tag{CE7}
\]
Thus $\bar P_d=1$ and $\bar P_4=1-\bar q\ge1/2$.
These exact low-degree identities are the physicality input needed
for the entry, rather than an estimate on a truncated coefficient.

\begin{theorem}[Weak fourth event gives the uniform entry]
For an actual canonical C1 source satisfying (CE3), the following
conditions are equivalent:
\[
 A_{\rm ev}/d\text{ is infinitesimal},\qquad
 \eta=f^2\text{ is infinitesimal},\qquad B\text{ is unbounded}.
                                                               \tag{CE8}
\]
Under any of these conditions the full original parent is tensor SOS
at every independent projective test and every real or complex state,
with no restriction on $e$.
\end{theorem}
\begin{proof}
The exact original source gives
\[
 \frac{A_{\rm ev}}d=\eta\frac{c_4}{c_d}
                  =\frac{\eta P_4}{2rP_d}.              \tag{CE9}
\]
The ratio $P_4/P_d$ is a positive bounded unit by (CE6)--(CE7).
If (CE9) is infinitesimal then $\eta/r$ is infinitesimal; since $r<1$,
$\eta$ is infinitesimal. Conversely, if $\eta$ is infinitesimal then
$r=1-a\eta$ has residue one because $a<2$, and (CE9) is infinitesimal.
The final equivalence follows from $B=\eta^{-1/2}$.
Now (CE1)--(CE3) give $k,f,g$ positive infinitesimals and the unchanged
positive bounded shapes and gaps of the uniform C1 source. The complete
uniform theorem, including its newly proved unbounded-$e$ branch,
applies for every positive $e=kf/g$. All quantities used for entry lie
in the coefficient field; no source/test rate comparison occurs.
\end{proof}

\paragraph{What the other weak templates mean on this same source row.}
The same calculation gives useful exact exclusions, not further SOS
claims. In (CE3),
\[
 \frac dc=\frac{2(1-r)P_d}{K(r+\eta)E},\qquad
 \frac{D_{\rm ev}}d=\frac{\mathcal B_5}{2J\eta P_d},\qquad
 \frac\alpha d=\frac{\kappa^2\zeta(r+\eta)P_\alpha}{2Q_cP_d}.
                                                               \tag{CE10}
\]
Since $1=r+a\eta\le2(r+\eta)$, the first ratio is bounded:
$c/d$ is never weak in this orientation. Also $P_\alpha,Q_c$ are positive
units, and $\kappa^2\zeta(r+\eta)=\kappa^2\zeta r+\kappa\zeta q$
is infinitesimal; hence $\alpha/d$ is always weak.
If $\eta$ is bounded, $q$ is infinitesimal and $\mathcal B_5$ is a unit,
so $D_{\rm ev}/d$ cannot be infinitesimal. If $\eta$ is unbounded,
$a=(1-r)/\eta$ is infinitesimal, $J$ is a unit, and (CE10) shows
$D_{\rm ev}/d$ infinitesimal. Thus
\[
       D_{\rm ev}/d\text{ weak}\quad\Longleftrightarrow\quad
       \eta\text{ unbounded}.                         \tag{CE11}
\]
The complete remaining endpoint ratio is
\[
 \frac\beta d=
 \frac{K(\kappa r+q)[1+(1-t)\zeta]-t\zeta P_d}
      {2t\zeta P_d}.                                  \tag{CE12}
\]
On the unbounded-$\eta$ branch, if $\beta/d$ is weak too, this forces
$q/(t\zeta)$ to have residue one. Equivalently $e^2/t$ has residue one;
in particular $q$ is infinitesimal. This is an exact necessary critical
balance, not an all-tests certificate for that branch.
For completeness,
\[
 \frac md=\frac{K(\kappa r+q)(1-t\zeta)(1+q\zeta)}
                  {2M\zeta P_d}.                      \tag{CE13}
\]
It records all weak-$m$ possibilities without comparing source scales.

\paragraph{Precise complement and a genuine entry obstruction.}
Within (CE3), the new natural criterion closes exactly the infinitesimal
$\eta$ slice. The bounded-unit and unbounded $\eta$ slices remain outside
the uniform theorem. The C1 stabilizer exchanges $A_{\rm ev},\alpha$,
$c,m$, and $D_{\rm ev},\beta$, but sends $k$ to $1/k$. Thus reflecting
one of these small-$k$ complements does not give a uniform small-$k$
entry. In particular the fact that $\alpha/d$ is always weak does not
close the row by a second use of (CE8).

Beyond (CE3), neither an ordinary actual Alice angle nor bounded
$\mathfrak T$ has been covered by this entry theorem. This distinction
cannot be removed by the weak-event hypothesis alone. Let $\varepsilon>0$
be infinitesimal and take the exact native source
\[
 k=1,\qquad f^2=g^2=\varepsilon,\qquad a=t=\tfrac12.
                                                               \tag{CE14}
\]
Every original retained event is positive: their leading values are
\[
 c=2,\quad m\sim\tfrac23,\quad D_{\rm ev}\sim\tfrac23,
 \quad\beta\sim2,\qquad
 d\sim4\varepsilon,\quad A_{\rm ev}\sim\varepsilon^2,
 \quad\alpha\sim\varepsilon^2.
\]
The displayed native contact has kernel and angle stationarity by the
six short native identities. On $Z=0$ the signed core is a sum of
squares, and its positive $d,c,A_{\rm ev}$ terms make the state
restriction positive definite, so the contact matrix has three positive
complementary eigenvalues. The inherited Laurent determinant has
nonnegative outer coefficients and two-dimensional support; its
logarithmic convexity makes this stationary zero a global minimum.
The full-parent inertia criterion, including the projective limits,
therefore makes this an actual physical source. These are exactly the
hypotheses of the previously proved stationary-contact sufficiency
argument, not a numerical positivity test. Its normalized signed coefficient is infinitesimal,
and $A_{\rm ev}/d,\alpha/d\sim\varepsilon/4$ are weak, so it belongs
to the weak-$A$ template. But $u_c=1$, fixed by the C1 stabilizer, and
$v_c=\varepsilon$. Neither orientation has the uniform requirement
$u_c$ infinitesimal. This is an obstruction to an exhaustive uniform
entry claim, not a counterexample to tensor SOS.

The resulting source frontier is therefore explicit: the weak-$A$
template is newly certified on the entire small-$u_c$/unbounded-
$\mathfrak T$ row, with arbitrary $e$; its other actual-angle rows,
and the unit/unbounded-$\eta$ portions of that row after other weak
orientations, still require additional analytical certificates or exits.
No archived all-tests label has been promoted to a proof here.
